\documentclass[11pt]{article}
\usepackage[margin=1in]{geometry}
\usepackage{amsmath,amssymb}
\usepackage{enumitem}
\usepackage{microtype}
\usepackage[hidelinks]{hyperref}

\setlength{\parindent}{0pt}
\setlength{\parskip}{4pt}

\newlist{qna}{enumerate}{1}
\setlist[qna]{leftmargin=1.6em,itemsep=3pt,topsep=4pt,label=\textbf{\arabic*.}}
\newcommand{\paper}[1]{\subsection*{#1}}

\newcommand{\q}[2]{\item \textbf{#1} #2}
\newcommand{\tiein}[1]{\par\smallskip\noindent\emph{Connection to the offline RL problem space and BCA:} #1\par\medskip}

\title{Literature Review: Offline Reinforcement Learning\\Version 2: Parts A and B}
\author{David Bayha\\ITSC 8110}
\date{November 2026}

\begin{document}
\maketitle



\section*{Part A: Research Papers}

Papers 1--8 carry over from Version 1; Papers 9--18 are new in Version 2. Every work is described with the same four themes, of which each paper names its three most central: distribution shift, the conservatism--improvement trade-off, uncertainty quantification and calibration, and theory vs.\ practice.



% ============================ PAPER 1 ============================
\paper{Paper 1: Conservative Q-Learning for Offline Reinforcement Learning (CQL)}
CQL trains the critic to be pessimistic about actions the dataset does not contain, so the policy is less drawn to accidentally inflated scores.
\begin{qna}

\q{What is the citation?}
Kumar, A., Zhou, A., Tucker, G., \& Levine, S. (2020). Conservative Q-learning for offline reinforcement learning. \emph{Advances in Neural Information Processing Systems, 33}, 1179--1191. Version reviewed: arXiv:2006.04779v3. \url{https://doi.org/10.48550/arXiv.2006.04779}

\q{What is the research area?}
The area is machine learning, specifically offline reinforcement learning (RL): learning a policy from a fixed log of past decisions with no further trial and error. CQL belongs to the ``pessimistic critic'' family, which lowers the scores of actions outside the data instead of restricting the policy.

\q{What is the research question?}
Can a critic trained to be systematically pessimistic about unseen actions produce a policy whose predicted value is a guaranteed floor on its true performance? Does that beat methods that keep the policy near the recorded behavior?

\q{What is the significance of the question?}
Offline RL is most valuable where live experimentation is expensive or dangerous but large logs already exist. Standard algorithms often fail there because the critic's guesses about unseen actions are never corrected and the policy is drawn to the inflated ones. A cheap, principled fix could let RL reuse logged data the way supervised learning reuses datasets.

\q{What are the contributions reported in this paper?}
The novel idea is an addition to the critic's objective that pushes down the scores of actions the critic likes too much and props up those of recorded actions. The paper proves, under stated assumptions and for large enough caution strength $\alpha$, that the learned value undercuts the true value and that, with high probability, the policy is at most a bounded amount worse than the logging policy (``safe improvement'').

\q{What is the research method?}
The paper combines algorithm design (two ready-to-use variants), mathematical analysis (mostly for simplified tabular cases, with extensions), and broad benchmarking against established methods on the D4RL suite and offline Atari.

\q{What are the results?}
On datasets recorded by a single behavior, CQL is close to the best prior method. On the mixed-source datasets, it usually wins by large margins (e.g., 91.1--110.5 on the three random-expert locomotion tasks, where competitors score at most 53.0). Only the CQL variants score above zero on the hardest mazes. Its predicted values sit below true returns on the three datasets checked, while the baselines compared there overestimate.

\q{What are the impacts and limitations of the contributions?}
CQL is a comparison baseline in many later papers reviewed here and showed pessimism can be added with a small code change. However, the caution knob $\alpha$ (or its auto-tuning threshold) was set differently across domains, and there is no established way to decide when to stop training without trying the policy live. One global $\alpha$ also applies the same strength to well-covered and poorly covered situations.

\q{How were the results evaluated or validated?}
The paper relies on mathematical proofs, D4RL scores averaged over four runs, direct comparisons of predicted values with actual returns, and ablations. The evidence is broad but statistically thin: tables of means without spreads, some baselines copied rather than rerun, and a few inconsistent tables.

\q{What are the ethical implications? How are they addressed?}
The reviewed version contains no broader-impact discussion. The ``safe improvement'' guarantee is only relative to the logging policy, probabilistic, and assumption-bound, so the word \emph{safe} could invite misplaced trust. A policy kept near the data also keeps the data's biases, and all validation is simulated.

\q{What themes are suggested by this research? (minimum 3)}
\emph{Distribution shift}: scores for unseen actions cannot be corrected offline. \emph{The conservatism--improvement trade-off}: pessimism must be dosed, here via $\alpha$. \emph{Theory vs.\ practice}: the guarantees are not proved for the deep networks actually used.

\end{qna}
\tiein{CQL exemplifies the ``one global pessimism knob'' paradigm. Applied to CQL, BCA (Bayesian Conformal Aggregation) keeps the critic-side idea but is designed to scale the conservative gap per sample by a calibrated width, on top of the single $\alpha$. CQL's appendix sketches a loosely related count-dependent penalty for the tabular case.}


% ============================ PAPER 2 ============================
\paper{Paper 2: A Minimalist Approach to Offline Reinforcement Learning (TD3+BC)}
TD3+BC shows that a standard online algorithm with one added ``stay near the recorded actions'' penalty matches more complicated offline methods on D4RL locomotion at a fraction of the compute.
\begin{qna}

\q{What is the citation?}
Fujimoto, S., \& Gu, S. S. (2021). A minimalist approach to offline reinforcement learning. \emph{Advances in Neural Information Processing Systems, 34}, 20132--20145. Version reviewed: arXiv:2106.06860v2. \url{https://doi.org/10.48550/arXiv.2106.06860}

\q{What is the research area?}
The area is offline RL for continuous control (robot-like tasks with smooth action spaces). TD3+BC belongs to the behavior-regularization family, which fights inflated scores for unseen actions by keeping the policy close to the recorded behavior.

\q{What is the research question?}
Can a standard online algorithm, changed as little as possible, reach state-of-the-art offline performance, and how much of the extra machinery in recent offline methods is necessary?

\q{What is the significance of the question?}
Every extra component is a choice that normally must be validated by trying policies, which requires the interaction offline RL is supposed to avoid. A minimal method helps isolate what works, raises the bar elaborate methods must clear, and stays cheap as datasets grow.

\q{What are the contributions reported in this paper?}
The paper audits two leading methods (CQL and Fisher-BRC), showing they owed much of their scores to poorly justified engineering tweaks, not just their stated ideas. It proposes TD3+BC: one penalty term in the actor's objective plus standardized inputs. It also demonstrates that the offline policies tested are unstable across episodes.

\q{What is the research method?}
The study combines critical analysis with experimentation. The actor maximizes $\lambda Q(s,\pi(s)) - (\pi(s)-a)^2$ averaged over recorded state-action pairs. The balance weight $\lambda$ is automatically rescaled by the typical size of the scores, so one setting is used for every dataset. Experiments cover fifteen standard locomotion datasets.

\q{What are the results?}
Summed over the fifteen tasks, TD3+BC essentially ties the best competitor (979.3 vs.\ 974.6) and beats its rerun of CQL (764.3). It trains in 39 minutes versus over four hours for CQL. TD3+BC does poorly on some datasets heavy in poor behavior (28.8 vs.\ CQL's 100.3 on a half-random, half-expert Walker2d mix) and on the hardest mazes.

\q{What are the impacts and limitations of the contributions?}
TD3+BC became what ReBRAC (Paper 4) calls ``a de-facto standard baseline to be compared against''. The paper offers no theory. A uniform pull toward recorded actions cannot tell good recordings from bad ones and can average over different behaviors. The diagnosed instability is left unsolved, and every sample receives the same constraint strength regardless of how well the data cover it.

\q{How were the results evaluated or validated?}
The paper averages scores over the final checkpoints rather than the luckiest one and reruns the two strongest competitors from their authors' code under one protocol. There are no statistical tests or confidence intervals and only five seeds despite high documented variance.

\q{What are the ethical implications? How are they addressed?}
A short statement notes that learning from logged data inherits its biases, warns of misuse, and argues the main effect is wider access. Arguably the paper's most useful safety contribution is documenting that offline policies can fail badly on individual episodes. No mitigations are proposed, and all evidence is simulated.

\q{What themes are suggested by this research? (minimum 3)}
\emph{Distribution shift}: inflated scores for unseen actions are fought by keeping the policy close to the recorded behavior. \emph{The conservatism--improvement trade-off}: TD3+BC uses one uniform pull toward the data and does poorly on some low-quality datasets and the hardest mazes. \emph{Theory vs.\ practice}: a few changed lines matched prior methods on locomotion tasks, and engineering, not just ideas, drove part of two prior methods' scores.

\end{qna}
\tiein{TD3+BC is the minimalist pole of the problem space, with one uniform constraint for every sample. Its successor (ReBRAC, Paper 4) found in its ablations that removing this constraint caused the largest overall drop. BCA is designed to strengthen this term per sample by a bounded, calibrated factor while keeping the global weight.}

% ============================ PAPER 3 ============================
\paper{Paper 3: Offline Reinforcement Learning with Implicit Q-Learning (IQL)}
IQL learns values using only the actions in the dataset, never asking ``what if'' about unseen actions, yet can still stitch together good behavior from suboptimal pieces.
\begin{qna}

\q{What is the citation?}
Kostrikov, I., Nair, A., \& Levine, S. (2022). Offline reinforcement learning with implicit Q-learning. In \emph{International Conference on Learning Representations (ICLR 2022)}. Version reviewed: arXiv:2110.06169v1. \url{https://doi.org/10.48550/arXiv.2110.06169}

\q{What is the research area?}
The area is offline RL. IQL takes a third path beside constraining the policy (TD3+BC) and penalizing the critic (CQL): \emph{in-sample} learning, in which value networks are evaluated only on recorded actions. The policy is extracted afterward by imitation weighted toward better-than-expected actions (AWR).

\q{What is the research question?}
Can an offline algorithm still plan ahead by chaining many steps of value backups toward optimal behavior while its value functions are trained only on logged actions, so that unseen actions are never scored?

\q{What is the significance of the question?}
Asking the critic about unseen actions is a main source of offline error, and the prior fixes the paper discusses compromise between safety and improvement. One-step in-sample alternatives avoid the problem but do poorly when fragments of good behavior must be stitched together. Removing the problem while keeping multi-step planning would ease the field's central tension, cheaply.

\q{What are the contributions reported in this paper?}
The novel core is ``implicit'' value learning. A state-value network trained with an asymmetric expectile loss, whose level acts as an optimism dial, estimates the best value achievable within the data. A second network then learns action values one step further, so the pair plans ahead without scoring an unseen action. The paper proves, for idealized exact solutions, that the dial's fully optimistic limit recovers the best in-data value.

\q{What is the research method?}
The paper combines algorithm design and a short limiting argument with benchmarks on standard locomotion, maze, hand-manipulation, and kitchen suites. It also runs a fine-tuning study in which offline training is followed by live interaction.

\q{What are the results?}
On locomotion, IQL roughly ties the best prior method (692.4 vs.\ 698.5). On the maze tasks that demand stitching it leads on five of six (378.0 vs.\ 303.6 total). Training takes about 20 minutes. After live fine-tuning IQL improves further (370 to 474 on mazes). IQL still loses individual tasks, for example to filtered imitation on two medium-expert locomotion datasets.

\q{What are the impacts and limitations of the contributions?}
IQL showed that in-sample learning need not give up long-horizon planning. However, continuous states usually carry only one recorded action, so the ``best in-data value'' relies on the network generalizing across similar states. Its fixed weighting also amplifies trustworthy and noisy advantages identically.

\q{How were the results evaluated or validated?}
The paper reports standard benchmark averages and a comparison corrected to use one dataset version for all methods. Baseline numbers come from mixed sources (one competitor's maze scores could not be reproduced and were copied). The result tables have no spreads, no procedure is stated for choosing per-domain settings, and only the optimism dial is ablated.

\q{What are the ethical implications? How are they addressed?}
The paper has no ethics discussion. Because the policy is built by reweighting recorded behavior, biases in the log are inherited and can be amplified whenever they scored well. The deliberate optimism of the upper expectile may be risky in safety-critical use, and the fine-tuning recipe involves live exploration that would carry physical risk on hardware.

\q{What themes are suggested by this research? (minimum 3)}
\emph{Distribution shift}: IQL sidesteps it during value learning, since training never queries unseen actions. \emph{The conservatism--improvement trade-off}: the expectile level sets how optimistic the values are, and the imitation weights' temperature sets how closely the policy follows the data. \emph{Theory vs.\ practice}: the proof concerns the fully optimistic limit of the dial, while in practice its two knobs are set per domain.

\end{qna}
\tiein{IQL is one of BCA's four host algorithms. Its imitation weight is the term BCA is designed to temper per sample with a calibrated width, while keeping IQL's global temperature.}

% ============================ PAPER 4 ============================
\paper{Paper 4: Revisiting the Minimalist Approach to Offline Reinforcement Learning (ReBRAC)}
ReBRAC shows that giving the simple TD3+BC baseline the quiet engineering upgrades of newer methods is enough to match or beat more complex ensemble-free algorithms on average.
\begin{qna}

\q{What is the citation?}
Tarasov, D., Kurenkov, V., Nikulin, A., \& Kolesnikov, S. (2023). Revisiting the minimalist approach to offline reinforcement learning. \emph{Advances in Neural Information Processing Systems, 36}, 11592--11620. Version reviewed: arXiv:2305.09836v2. \url{https://doi.org/10.48550/arXiv.2305.09836}

\q{What is the research area?}
The area is offline RL, in the behavior-regularization lineage (BRAC, TD3+BC). The paper deliberately proposes no new principle.

\q{What is the research question?}
How far can a minimalist baseline be improved purely by adopting the small design choices that recent methods shipped quietly, and which of them account for the gain? The choices are deeper networks, layer normalization, bigger batches, separate penalty weights for actor and critic, and a longer planning horizon.

\q{What is the significance of the question?}
If gains credited to new algorithms actually come from quiet engineering, the field mis-attributes progress and practitioners cannot tell what is worth adopting. Establishing how strong a cheap baseline can become raises the bar new methods should clear and gives practitioners a fast default.

\q{What are the contributions reported in this paper?}
The paper audits which of five design choices seven recent algorithms adopt and introduces ReBRAC: TD3+BC plus five upgrades, including separate ``stay near the data'' penalties on the actor and on the critic. It evaluates ReBRAC on 51 datasets plus an offline-to-online study, runs large-scale ablations, and supplies code and per-dataset settings.

\q{What is the research method?}
The method is design by composition, then large-scale benchmarking and ablation, with no new theory. Diagnostic runs remove or add one upgrade at a time.

\q{What are the results?}
On average, ReBRAC is the best of the compared ensemble-free methods on mazes (76.8) and hand-manipulation tasks (58.6) and a close second on locomotion (81.2 vs.\ SAC-RND's 82.6). It fine-tunes online from 55.0 to 91.1. In the ablations, removing the actor's stay-near-the-data penalty makes the method collapse (72.2 to 11.4), removing critic layer normalization halves the score, and the critic-side penalty adds almost nothing on average.

\q{What are the impacts and limitations of the contributions?}
ReBRAC supplied a stronger cheap baseline (Papers 6 and 8 here compare against it) and practical guidance from its ablations. However, its two penalty weights are still re-tuned per dataset by trying policies live (best values span two orders of magnitude), and constraint strength is uniform across samples \emph{within} a dataset. Most tasks recorded by humans remain unsolved, and everything is simulated.

\q{How were the results evaluated or validated?}
The evaluation is fairly rigorous, with separate tuning and reporting seeds, final-checkpoint reporting, some baselines re-tuned under the same protocol, and a control showing the gains come from the design choices rather than implementation differences. Still, model selection uses real environment returns, which a true offline user lacks, and copied baselines come from other papers' protocols, which may be more lenient.

\q{What are the ethical implications? How are they addressed?}
Ethics are not discussed beyond scientific scope. Implicitly, staying near logged behavior inherits its biases and unsafe habits, and tuning by live trial runs contradicts the offline premise in risky domains. Performance also dipped on one task during live fine-tuning, which outside simulation would mean worse behavior in operation.

\q{What themes are suggested by this research? (minimum 3)}
\emph{Distribution shift}: the safeguard against unsupported actions is the actor's stay-near-the-data penalty, and the method collapses without it. \emph{The conservatism--improvement trade-off}: constraint strength is uniform across samples within a dataset. \emph{Theory vs.\ practice}: incidental choices may account for much reported ``algorithmic'' advance.

\end{qna}
\tiein{ReBRAC's ablations identify the actor's stay-near-the-data term as the lever that matters most, but its strength remains a single scalar re-tuned per dataset by live evaluation. BCA's calibrated weighting targets the same lever, dosed per sample within a bounded range on top of the host's tuned scalar.}


% ============================ PAPER 5 ============================
\paper{Paper 5: Importance Weighted Actor-Critic for Optimal Conservative Offline Reinforcement Learning (A-Crab)}
A-Crab proves that weighting each sample's error in the critic's penalty, instead of squaring it, turns a slow-but-practical method into one that attains the optimal statistical rate under its assumptions.
\begin{qna}

\q{What is the citation?}
Zhu, H., Rashidinejad, P., \& Jiao, J. (2023). Importance weighted actor-critic for optimal conservative offline reinforcement learning. \emph{Advances in Neural Information Processing Systems, 36}, 49579--49602. Version reviewed: arXiv:2301.12714v2. \url{https://doi.org/10.48550/arXiv.2301.12714}

\q{What is the research area?}
The area is theoretical offline RL with general function approximation. The paper connects adversarial actor-critic methods (ATAC) with importance weighting, which corrects for how rarely some behavior appears in the data.

\q{What is the research question?}
Can one practical algorithm achieve the statistically optimal learning speed (error shrinking like $1/\sqrt{N}$ as the data grow) under a weak, averaged assumption about how well the data cover alternative behaviors? Prior practical methods had slower proven rates or needed stronger, worst-case coverage.

\q{What is the significance of the question?}
A guarantee is only as useful as its assumptions. Real logs rarely cover every possible policy, and worst-case coverage assumptions collapse when a single rarely logged action makes a likelihood ratio enormous.

\q{What are the contributions reported in this paper?}
The key step replaces the squared Bellman-error penalty with the largest \emph{importance-weighted average} error (a Bellman error is the critic's disagreement with its own one-step lookahead). The paper adds an optimal-rate proof that drops the strong ``completeness'' closure assumption predecessors needed, and a weaker coverage measure. It also proves the squared penalty fails to reach the optimal rate on some problems and provides a practical variant with public code.

\q{What is the research method?}
Mathematics comes first: the paper analyzes a game in which the actor tries to beat the logged behavior while the critic minimizes that improvement plus the weighted error penalty. A tractable approximation then restricts the weights to a boxed range, yielding a closed-form penalty, and is benchmarked against ATAC on twelve simulated locomotion datasets.

\q{What are the results?}
Under its assumptions, A-Crab attains the optimal $1/\sqrt{N}$ rate against any comparator policy the data cover and guarantees a policy no worse than the logging policy, up to a vanishing error, across a wide range of penalty weights. In practice it usually trains more stably than ATAC and scores better on several datasets (about 110 vs.\ a noisy 90--100 on one task), though results appear only as training curves.

\q{What are the impacts and limitations of the contributions?}
The main impact is conceptual: in the paper's analysis, the \emph{form} of the error penalty is what separates a tractable-but-slow method from a rate-optimal one. However, the assumptions (finite action sets, idealized actor updates, well-behaved weights) are far from deep-learning practice. The per-sample weighting also reduces in practice to a closed-form penalty with a single grid-searched constant per dataset.

\q{How were the results evaluated or validated?}
The theorems have full appendix proofs. The empirical side is thin: it shows only training curves against a single baseline, with no final-score table or statistical tests. No results are shown for the new penalty without the squared one, and the theoretical rate improvement is not checked numerically.

\q{What are the ethical implications? How are they addressed?}
There is no ethics discussion. Optimal-rate language could create false confidence for clinical or robotic use. The guarantees rest on assumptions that are hard to verify from logged data, concern average return rather than per-state safety, and do not cover the benchmarked version. Because the method competes only with behaviors the data cover, logging bias can carry over into the learned policy.

\q{What themes are suggested by this research? (minimum 3)}
\emph{Distribution shift}: guarantees are sought under partial coverage, with optimal rates against covered behaviors only. \emph{The conservatism--improvement trade-off}: the analysis guarantees a policy no worse than the logging policy, up to a vanishing error, across many penalty weights. \emph{Theory vs.\ practice}: the tested algorithm is not the analyzed one.

\end{qna}
\tiein{A-Crab's per-sample reweighting of errors is a loose theoretical cousin of BCA's calibrated widths. The paper proves its importance weighting attains the optimal rate under its assumptions, a result that does not carry over to BCA. Its practical reduction to one searched constant per dataset is the kind of gap BCA is designed to target.}

% ============================ PAPER 6 ============================
\paper{Paper 6: A Clean Slate for Offline RL (Unifloral)}
This paper rebuilds twelve existing algorithms in one fast codebase and grades ten of them by how well each scores for a given amount of live testing spent on tuning.
\begin{qna}

\q{What is the citation?}
Jackson, M. T., Berdica, U., Liesen, J., Whiteson, S., \& Foerster, J. N. (2025). A clean slate for offline reinforcement learning. \emph{Advances in Neural Information Processing Systems, 38}. Version reviewed: arXiv:2504.11453v1. \url{https://doi.org/10.48550/arXiv.2504.11453}

\q{What is the research area?}
The area is offline RL methodology, evaluation, and software rather than a new learning rule. Its problem family is hyperparameter and policy selection under limited online access, together with fair benchmarking.

\q{What is the research question?}
How should offline RL methods be defined, implemented, and compared so that reported scores honestly reflect the live interaction spent tuning them? Can one cleanly implemented design space, containing prior algorithms as settings, be used to discover better ones?

\q{What is the significance of the question?}
Offline RL's selling point, learning without touching the real world, is undermined if published scores quietly depend on many live test episodes used to pick settings. Numbers from different papers then mix algorithmic quality with tuning effort.

\q{What are the contributions reported in this paper?}
The paper distinguishes four offline RL settings by when live interaction is allowed. Its evaluation protocol simulates tuning, so performance is plotted against the live-testing budget. It draws a family tree of prior methods, reimplements twelve of them, and builds Unifloral, one configurable implementation covering most of the tree (it omits CQL). It creates two new algorithms (TD3-AWR and MoBRAC) by mixing settings, with no new code.

\q{What is the research method?}
The method is conceptual analysis, protocol and framework construction, then benchmarking, with no new theory. For each algorithm--dataset pair, 20 policies are trained with settings sampled from the original papers' ranges and their episode returns stored. Then 500 simulated tuning sessions select among them under a stated budget.

\q{What are the results?}
None of the ten prior methods tested dominates. ReBRAC leads on five of nine datasets at some budget and IQL on four. The reimplementations are faster than earlier libraries, in some cases dramatically. TD3-AWR beats each of its parent methods on most datasets (ReBRAC on six of nine, IQL on seven). With only a few noisy test episodes, selection can be fooled by erratic ``distractor'' policies (shown for ReBRAC on one dataset).

\q{What are the impacts and limitations of the contributions?}
The paper offers the field a common yardstick and toolkit: hidden tuning becomes an explicit, plotted cost. However, only the pre-deployment setting is implemented, rankings depend on the original papers' tuning ranges, and everything uses one benchmark suite. Some fixed configuration choices also appear to degrade certain algorithms relative to their originals, and the two headline average speedups are not recoverable from the reported data.

\q{How were the results evaluated or validated?}
The evaluation uses stored-episode tuning simulations with confidence intervals, a locomotion reproduction table said to match an earlier library's scores (reference numbers not shown), and a distractor-policy analysis. There are no formal statistical tests, and the intervals reflect selection randomness over one pool of trained policies rather than retraining. No reported check confirms the unified implementation matches each standalone file.

\q{What are the ethical implications? How are they addressed?}
There is no dedicated ethics section, but the paper can be read as a research-integrity intervention: undisclosed live tuning can inflate deployability claims, and pricing that budget makes the inflation visible. It leaves unaddressed that selecting policies by average return ignores worst-case outcomes in the safety-critical domains it cites, and that distractor policies are the kind a careless real selection process could deploy. Every benchmark is also simulated.

\q{What themes are suggested by this research? (minimum 3)}
\emph{The conservatism--improvement trade-off}: the strongest methods in the paper's tests still tie the policy to the data. \emph{Uncertainty quantification and calibration}: policy selection is itself uncertain, and the main budget curves have 95\% confidence intervals over simulated tuning sessions. \emph{Theory vs.\ practice}: live tuning, which offline RL in theory avoids, is priced in episodes, with budget curves instead of headline numbers.

\end{qna}
\tiein{Unifloral is the baseline reference bundled with BCA's code, and its budget-aware protocol is a suitable measuring stick for BCA's aim. If calibrated per-sample caution reduced the search over global knobs, the payoff would appear as better scores at small live-testing budgets. That remains untested.}

% ============================ PAPER 7 ============================
\paper{Paper 7: Quantile Q-Learning: Revisiting Offline Extreme Q-Learning with Quantile Regression (QQL)}
QQL replaces a sensitive hand-tuned knob in a recent offline algorithm, a single ``optimism temperature'' per dataset, with a per-state scale learned from the data.
\begin{qna}

\q{What is the citation?}
Gao, X., Li, S., Cai, Y., \& Yu, W. (2026). Quantile Q-learning: Revisiting offline extreme Q-learning with quantile regression. \emph{Transactions on Machine Learning Research}. Version reviewed: arXiv:2511.11973v2. \url{https://doi.org/10.48550/arXiv.2511.11973}

\q{What is the research area?}
The paper is in offline RL, downstream of the in-sample family (IQL, XQL), which learns values only on recorded actions. It addresses temperature selection, or how optimistic the value estimates should be, by borrowing quantile regression (learning specific percentiles rather than averages) from distributional RL.

\q{What is the research question?}
Can an XQL-style algorithm stay stable and competitive across many domains with one fixed configuration if its temperature is estimated from the data per state instead of tuned per dataset?

\q{What is the significance of the question?}
Per-dataset tuning normally means trying candidate policies live, the interaction offline RL is meant to avoid. XQL's temperature, with reproduced values spanning 0.5 to 10 across datasets, feeds both value learning and the policy's imitation weights, so a bad choice can destabilize both. A learned per-state scale would remove that knob.

\q{What are the contributions reported in this paper?}
The paper shows that, under an assumed noise model, the temperature equals a constant times the gap between two specific percentiles of the action-value distribution, so quantile regression can learn it without hand-tuning. It builds and benchmarks QQL, which combines two percentile value networks, a bias-corrected value target, a policy update whose optimism varies by state, and a cautious generalization term for unseen actions.

\q{What is the research method?}
The paper derives its result under stated assumptions, then designs an algorithm and tests it on 19 datasets from D4RL and NeoRL2. All settings are fixed across datasets, while baselines are allowed their tuned values.

\q{What are the results?}
With one fixed configuration, QQL posts the best average on 10 of 17 standard datasets. When XQL must also share one temperature, QQL wins clearly on locomotion (85.6 vs.\ 73.6) and mazes. Its training curves are smoother than a stabilized XQL variant's on the four datasets plotted. It is weak on both Hammer manipulation datasets and scores below the dataset's own average on the one safety-constrained task.

\q{What are the impacts and limitations of the contributions?}
QQL is evidence that XQL's temperature can be estimated per state from data. The assumed noise model appears to fit the reported residuals only loosely; they look heavier-tailed on the four datasets analyzed. The two learned percentiles can cross, which is patched by hand. Some ablation claims overstate the figures.

\q{How were the results evaluated or validated?}
The evaluation uses five-run averages (three-run on NeoRL2), ablations, knob sweeps, a synthetic test, and an error-distribution analysis. The benchmark scores have no confidence intervals or significance tests, and baselines come from mixed sources. The choice of shared temperature for the head-to-head comparison is unexplained, and the synthetic test injects exactly the noise it then detects.

\q{What are the ethical implications? How are they addressed?}
The paper discusses none. Its claim that fixed settings make the method practical for consequential use rests entirely on simulation, and constraint violations on the safety-constrained task are not reported. A misestimated scale in sparsely covered states can mean \emph{less} caution where mistakes may cost most.

\q{What themes are suggested by this research? (minimum 3)}
\emph{The conservatism--improvement trade-off}: optimism is set per state instead of by one temperature, and wider value spread across recorded actions triggers stronger caution. \emph{Uncertainty quantification and calibration}: the learned scale is a modeled quantity with no coverage guarantee. \emph{Theory vs.\ practice}: hyperparameter sensitivity is arguably relocated, not removed, into a quantity learned under an assumed noise model; its two controls look insensitive on the three datasets swept.

\end{qna}
\tiein{QQL is a close non-conformal precedent for BCA: its per-state scale acts on the same imitation weight that BCA modulates in its IQL host. QQL's scale is learned under a parametric noise model with no guarantee. BCA's conformally calibrated width is designed as an alternative, though its coverage concerns held-out residuals, not the policy.}

% ============================ PAPER 8 ============================
\paper{Paper 8: UNIQ: Conformal Calibration for Adaptive Conservatism in Offline Reinforcement Learning (UNIQ)}
UNIQ gives a standard offline learner (IQL) a low-memory, per-state caution dial driven by conformally rescaled ensemble disagreement.
\begin{qna}

\q{What is the citation?}
Upadhyay, A. (2026, July 11). \emph{UNIQ: Conformal calibration for adaptive conservatism in offline reinforcement learning} [Poster presentation]. ICML 2026 Workshop on Decision-Making from Offline Datasets to Online Adaptation: Black-Box Optimization to Reinforcement Learning, Seoul, South Korea. Version reviewed: arXiv:2606.07592v1. \url{https://doi.org/10.48550/arXiv.2606.07592}

\q{What is the research area?}
The paper is in offline RL and builds on IQL. The problem family is pessimism under distribution shift: how much to distrust value estimates in poorly covered situations. Unlike global-knob methods (CQL, IQL, TD3+BC) and large-ensemble methods (SAC-N, EDAC; Paper 17), UNIQ pairs a small ensemble of value heads with conformal calibration.

\q{What is the research question?}
Can IQL's global optimism setting be replaced by a per-state setting driven by conformally normalized ensemble disagreement, with tighter caution where data are sparse and looser where they are dense, without the memory cost of large ensembles?

\q{What is the significance of the question?}
The best setting of most mainstream caution knobs changes per dataset and is normally found by trying policies live, an option unavailable offline. Raw ensemble disagreement has no common scale across tasks, so fixed penalties transfer poorly.

\q{What are the contributions reported in this paper?}
The main contribution is a three-part mechanism. A small ensemble of value heads estimates uncertainty from its spread. A conformal step measures held-out prediction errors and converts raw spread into a unit-free ratio. A controller maps that ratio to a state-specific optimism setting, with a matching penalty in the value targets. Supporting contributions include a formal appendix and a benchmark comparison that keeps the losing rows.

\q{What is the research method?}
The method is algorithm design with light formal analysis and benchmarking. UNIQ uses three heads per optimism level, and held-out transitions periodically recalibrate the error scale. Two hand-set configurations cover nine simulated locomotion datasets.

\q{What are the results?}
UNIQ averages 85.2 across the nine tasks versus a copied IQL figure of 81.6 and ranks third overall, with measured memory of 250 MB against an estimated 2.5 GB for the ensemble leader. In the single-seed, four-task ablations, the hand-set penalty coefficient appears to matter more at matched settings than calibration or adaptivity, and a single head per level, which gives no uncertainty signal, nearly matches the full method (59.1 vs.\ 59.4).

\q{What are the impacts and limitations of the contributions?}
UNIQ is a proof of concept that conformal calibration can be added cheaply to a standard learner. The conformal guarantee is marginal: it covers held-out Bellman residuals under one task-wide threshold, not individual states, and holds only at each recalibration of a moving ensemble. The gains are partly confounded with per-task penalty settings, the evidence is locomotion-only, and several numbers conflict across tables, text, and appendix.

\q{How were the results evaluated or validated?}
The evaluation compares three-seed benchmark scores with eight baselines copied (not rerun) from a prior report. Realized coverage and the distribution of the uncertainty signal are never reported, so the headline 3.6-point gain over IQL is plausible but weakly established, and the IQL comparison is not matched.

\q{What are the ethical implications? How are they addressed?}
The paper has no ethics discussion. Its language of distribution-free finite-sample guarantees could imply a safety assurance the method does not provide. The coverage concerns one-step Bellman residuals on held-out logged transitions, not the states the learned policy will visit. The controller also grows \emph{more} optimistic wherever ensemble heads agree, which is dangerous if they share the same error.

\q{What themes are suggested by this research? (minimum 3)}
\emph{The conservatism--improvement trade-off}: the paper's premise is that one global setting is too careful in dense regions and may be too lax in sparse ones. \emph{Uncertainty quantification and calibration}: the paper argues that raw disagreement needs a common scale. \emph{Theory vs.\ practice}: conformal coverage on residuals is not shown to carry over to the derived signal or to states the policy visits.

\end{qna}
\tiein{Of the papers reviewed here, UNIQ is closest to BCA: both bring conformal ideas to prediction-error-driven caution. UNIQ's guarantee does not formally extend to its uncertainty signal (rescaled ensemble spread), whereas BCA's signal is the conformally calibrated width of the errors themselves. BCA's guarantee is still marginal and on residuals, not per sample or on the policy.}

% ============================ PAPER 9 ============================
\paper{Paper 9: Bayesian Flow Networks for Offline Trajectory Planning (BFN-RL)}
BFN-RL swaps the diffusion model inside a trajectory planner for a Bayesian Flow Network, so one kind of generator can plan over both grid-world (discrete) and robot-control (continuous) states.
\begin{qna}

\q{What is the citation?}
Killingberg, L., \& Langseth, H. (2026). \emph{Bayesian flow networks for offline trajectory planning} (arXiv:2608.25163v1) [Preprint]. arXiv. \url{https://doi.org/10.48550/arXiv.2608.25163}

\q{What is the research area?}
The area is offline RL treated as sequence generation rather than value learning. In the line of Diffuser and the Decision Diffuser, the planner generates a multi-step future state sequence conditioned on a desired return or goal. It then recovers actions with a learned inverse-dynamics model, a network that infers which action leads from one state to the next.

\q{What is the research question?}
Can a Bayesian Flow Network (BFN) serve as a single generator for offline trajectory planning that handles discrete and continuous states natively while staying competitive with diffusion-based planners in each setting?

\q{What is the significance of the question?}
Many planning problems are discrete or mix discrete and continuous variables, yet diffusion planners are built for continuous vectors and need a separate categorical formulation for discrete states. One generator that covers both would simplify planner design, and a matched comparison can indicate whether the choice of generator deserves attention at all.

\q{What are the contributions reported in this paper?}
The paper contributes BFN-RL, a planner whose generator refines the parameters of a distribution over a state trajectory instead of denoising a noisy one, so the same machinery applies to discrete and continuous spaces. Its discrete planner pairs a categorical BFN with a learned inverse-dynamics model, using an inpainting rule that pins known states during sampling. A cross-domain study covers D4RL locomotion and five discrete tasks.

\q{What is the research method?}
This is an engineering study that designs and benchmarks a planner, with no theoretical analysis. The discrete testbed is three MiniGrid tasks, slippery FrozenLake, and two-box Sokoban, where the BFN and diffusion generators share data, network, conditioning, and seeds. For continuous control, a Gaussian BFN is scored on nine D4RL locomotion datasets.

\q{What are the results?}
In the matched discrete comparison, neither generator dominates. Four of the five tasks are close in success rate, and BFN-RL wins clearly only on Sokoban (68.7\% vs.\ 59.0\%). On D4RL it averages 77.0, behind the Decision Diffuser (81.8) and level with IQL (77). It falls below plain behavior cloning on walker2d-medium (66.9 vs.\ 75.3).

\q{What are the impacts and limitations of the contributions?}
The paper is a proof of concept that BFNs can drive offline planners, with a matched-baseline protocol designed to isolate the generator's effect. By the authors' own account, BFN-RL does not beat diffusion across the board. No component is ablated, per-task settings are reported without explanation of how they were chosen, and the discrete datasets contain only successful trajectories.

\q{How were the results evaluated or validated?}
The discrete tasks use a matched protocol with three training seeds. Continuous control uses D4RL normalized scores over three seeds, with every baseline number copied from the Decision Diffuser paper. The matched design is the strongest part. Confidence is limited by the three seeds and the lack of significance tests.

\q{What are the ethical implications? How are they addressed?}
The paper has no ethics statement or limitations section. Its introduction motivates offline learning with self-driving cars and clinical treatment, where exploration is dangerous. Yet a generative planner can propose state sequences the system cannot reach, and conditioning on returns outside the data may push generation into unfamiliar regions with no uncertainty signal. Everything is simulated.

\q{What themes are suggested by this research? (minimum 3)}
\emph{Distribution shift}: with no critic, the inflated-value problem does not arise in the same form. \emph{The conservatism--improvement trade-off}: the planner seeks improvement by return or goal conditioning, yet on walker2d-medium its three-seed mean falls below the copied behavior-cloning score. \emph{Theory vs.\ practice}: a matched baseline is designed to isolate the generator, but copied baselines and three seeds weaken the continuous comparison.

\end{qna}
\tiein{BFN-RL represents the critic-free branch of the problem space. With no critic there are no Bellman residuals (the Bellman errors of Paper 5), so BCA's calibration cannot attach directly. One loose analogue is BFN-RL's fixed sampling temperature, which arguably plays the role of a global caution knob.}

% ============================ PAPER 10 ============================
\paper{Paper 10: Conformal Prediction Beyond the Horizon: Distribution-Free Inference for Policy Evaluation (Conformal Policy Evaluation)}
This paper adds a conformal calibration layer to a distributional RL model, so a policy's long-run return comes with a prediction interval checked against held-out data.
\begin{qna}

\q{What is the citation?}
Gan, F., Lu, Y., Zhang, Y., \& Liu, Y. (2025). Conformal prediction beyond the horizon: Distribution-free inference for policy evaluation. \emph{Advances in Neural Information Processing Systems, 38}. Version reviewed: arXiv:2510.26026v1. \url{https://doi.org/10.48550/arXiv.2510.26026}

\q{What is the research area?}
The paper applies statistics to RL. It uses conformal prediction (prediction intervals that assume no particular form for the data's distribution) for policy evaluation and builds on distributional RL.

\q{What is the research question?}
How can prediction intervals for a target policy's infinite-horizon discounted return be built from finite logged trajectories to reach a chosen coverage level without depending on a correctly specified return model, both on-policy and off-policy?

\q{What is the significance of the question?}
An expected return hides the spread of outcomes: a policy can look good on average and still do badly from some states. Distributional RL's quantiles can undercover (miss the true return more often than claimed). Conformal prediction offers model-agnostic validity, but its exchangeability assumption (that data order does not matter) fails for sequential data, and infinite-horizon returns are never fully observed.

\q{What are the contributions reported in this paper?}
The paper reports five contributions. It builds a pseudo-return label from $k$ observed discounted rewards plus a tail sampled from the learned return distribution. It subsamples calibration transitions to weaken time dependence, and uses weighted resampling to correct for state shift and, off-policy, action shift. It gives asymptotic coverage bounds and a simulation study with public code.

\q{What is the research method?}
The paper combines algorithm design, asymptotic analysis, and simulation. Each calibration transition is scored by the gap between its pseudo-return and the value predicted by a return-distribution model fitted on training transitions. The simulations include a two-state problem, a two-dimensional continuous problem, and Mountain Car.

\q{What are the results?}
In theory, coverage in the large-sample limit is at least the target minus weight-estimation error and a model-error term that shrinks geometrically in $k$. In the two main simulations, conformal coverage stays near the 90\% target (medians roughly 0.89--0.92) while raw quantile intervals undercover (roughly 0.84--0.87), though conformal intervals are wider in the continuous example. On Mountain Car, on-policy coverage is near 0.90 but off-policy coverage stays near 0.85.

\q{What are the impacts and limitations of the contributions?}
The paper shows that a conformal layer can sit on top of an ordinary distributional RL estimator. But the guarantees are asymptotic and hold on average over test states. The authors call the choice of $k$ underexplored, and two-state coverage rises with it, from 0.87 at $k=1$ to 0.94 at $k=8$. A second tuning parameter was set using long simulated rollouts that are unavailable offline, and the experiments are small and synthetic.

\q{How were the results evaluated or validated?}
The theory rests on appendix proofs under Markov, stationarity, bounded-density, and overlap assumptions. In the simulations, long rollouts reveal the true return, so coverage is measured directly. At least one claim goes beyond the figures: the text calls Mountain Car coverage close to the target, while its off-policy medians sit near 0.85.

\q{What are the ethical implications? How are they addressed?}
There is no ethics section, and the authors answer no to the NeurIPS checklist's broader-impact and safeguards questions, although the motivation is clinical treatment and autonomous systems. An interval labeled 90\% can create unwarranted confidence when its assumptions fail, and because coverage is only an average over states, a patient subgroup can be undercovered. The paper discusses neither risk.

\q{What themes are suggested by this research? (minimum 3)}
\emph{Distribution shift}: weighted resampling is meant to correct for state shift and, off-policy, action shift, but off-policy coverage on Mountain Car stays near 0.85. \emph{Uncertainty quantification and calibration}: raw distributional quantiles undercover in the simulations, and held-out calibration narrows or closes the gap. \emph{Theory vs.\ practice}: coverage survives an imperfect model only approximately and asymptotically, and both tuning parameters were assessed against ground-truth returns unavailable in logged data.

\end{qna}
\tiein{This paper is an evaluation-side counterpart of BCA: both calibrate bootstrapped quantities on held-out transitions. Its theory is a warning for BCA: with one-step targets, coverage of the true return is certified only up to the model's own error. Its main experiments use constant-width intervals, whereas BCA's fitted scale model lets the width vary per sample.}

% ============================ PAPER 11 ============================
\paper{Paper 11: Bellman Calibration for V-Learning in Offline Reinforcement Learning (Bellman Calibration)}
Bellman calibration checks whether a trained value function's numbers are systematically mis-scaled and, if so, adjusts them with a one-dimensional correction learned from logged transitions.
\begin{qna}

\q{What is the citation?}
van der Laan, L., \& Kallus, N. (2026). \emph{Bellman calibration for V-learning in offline reinforcement learning} (arXiv:2512.23694v2) [Preprint]. arXiv. \url{https://doi.org/10.48550/arXiv.2512.23694}

\q{What is the research area?}
The paper sits between statistics and RL. It studies post-hoc calibration of a trained value function for a fixed target policy, the policy-evaluation side of offline RL.

\q{What is the research question?}
Can a trained value predictor be checked for systematic errors in its numerical scale and repaired by a single scalar transformation of its outputs, learned from logged transitions?

\q{What is the significance of the question?}
Value estimates feed real decisions, so a mis-scaled value function can distort comparisons between policies even when its ranking of states is sensible. Existing guarantees for fitted value learning typically rest on strong structural conditions such as Bellman completeness, and the authors find much less model-agnostic theory for diagnosing and repairing scale after training.

\q{What are the contributions reported in this paper?}
The paper reports four contributions. It defines Bellman calibration: among states that share a predicted value, the average one-step backup should equal that prediction. It gives a calibration-error measure with a doubly robust off-policy estimator (valid if either helper model is right), and a bound tying value error to calibration error plus what no rescaling can fix. Its post-hoc procedure, Iterated Bellman Calibration, has finite-sample bounds assuming neither completeness nor a correctly specified value class.

\q{What is the research method?}
The paper is mainly theory and method, supported by simulations and a benchmark audit. The calibrator repeatedly regresses one-step backup targets on the original predictions, so every iterate stays a transformation of the initial value function. The audit calibrates a frozen neural critic on the Deep OPE benchmark's Hopper data.

\q{What are the results?}
In theory, calibration error is bounded, up to helper-model and iteration terms, by one-dimensional statistical rates. In simulation, corrections cut value error where the scale was wrong: about a fifth under early stopping but under a tenth with the saturating reward. On Hopper, the hybrid calibrator cuts calibration error to about 0.78 of its raw level for the better-trained of two critics, though final policy-value error barely moves (0.978 of its raw level).

\q{What are the impacts and limitations of the contributions?}
The paper offers a cheap, model-agnostic audit-and-repair step for any critic and explains when rescaling can help (good ranking, wrong scale) and when it cannot. The authors stress that calibration is not accuracy, and the fully iterated isotonic version lacks a value-error guarantee. The simulations use only clipped importance weights, a special case of the doubly robust targets the theory centers on.

\q{How were the results evaluated or validated?}
The results are validated by appendix proofs, simulations scored against ground truth, and a Hopper audit with held-out scoring. But the simulated failure modes are the ones the method is designed to fix, and several settings go unreported. The mis-scaled-reward value-error ratios differ without explanation between a figure and a table (0.42 vs.\ 0.73--0.75), and no external recalibration or evaluation method is compared.

\q{What are the ethical implications? How are they addressed?}
None are discussed. Because calibration conditions only on the predicted value, errors for one group can cancel against opposite errors elsewhere at the same prediction level, so a low calibration error could hide systematic harm to a subpopulation. A small calibration error might also be over-read as validation, and all empirical evidence comes from simulators.

\q{What themes are suggested by this research? (minimum 3)}
\emph{Distribution shift}: a target policy's Bellman backups are estimated from another policy's data. \emph{Uncertainty quantification and calibration}: calibration is carried over to long-horizon value predictions by a post-hoc fix; it concerns the reliability of the mean, not interval coverage. \emph{Theory vs.\ practice}: on the better-trained Hopper critic, calibrators with value-error guarantees cut calibration error while those without raised it, though elsewhere both kinds helped.

\end{qna}
\tiein{Bellman calibration complements BCA: BCA reads uncertainty from the size of Bellman residuals, while this paper targets their average bias among states with similar predicted values. Its calibration error could serve as a companion diagnostic for BCA's host critics, though its guarantees cover a frozen critic and a fixed policy, not actor-critic training.}

% ============================ PAPER 12 ============================
\paper{Paper 12: Weighted Bayesian Conformal Prediction (WBCP)}
WBCP gives a conformal threshold its own error bar: alongside each prediction interval it reports how firmly the calibration data support that interval's width.
\begin{qna}

\q{What is the citation?}
Lou, X., \& Luo, P. (2026). \emph{Weighted Bayesian conformal prediction} (arXiv:2604.06464v2) [Preprint]. arXiv. \url{https://doi.org/10.48550/arXiv.2604.06464}

\q{What is the research area?}
The area is statistical machine learning rather than RL. Within conformal prediction, the paper sits where weighted conformal methods for covariate shift (test inputs distributed differently from the calibration inputs) meet Bayesian quadrature conformal prediction (BQ-CP, Paper 18).

\q{What is the research question?}
Can the posterior over the conformal threshold that BQ-CP provides for i.i.d.\ data be extended to weighted conformal prediction? The aim is one procedure that stays valid under covariate shift or local weighting and also reports how reliable its own threshold is.

\q{What is the significance of the question?}
A conformal threshold is a single number that says nothing about how firmly the calibration data pin it down. That matters wherever interval width drives decisions and calibration data are unevenly spread. The authors argue that the two existing remedies, BQ-CP's threshold posterior and weighted conformal prediction, had not been combined.

\q{What are the contributions reported in this paper?}
The paper reports four contributions. WBCP replaces BQ-CP's flat Dirichlet model with a weighted one scaled by Kish's effective sample size (how many equally weighted points the weighted sample is worth). Four theoretical results, given as proof sketches, address why this scaling is variance-matching and when BQ-CP's guarantees carry over. The paper also gives a spatial version with an adaptive bandwidth (AdaGeoBCP) and an empirical study.

\q{What is the research method?}
The paper combines algorithm design and asymptotic proof sketches with benchmarking and ablations. Many Monte Carlo draws from the weighted Dirichlet each give a threshold; a high quantile of those thresholds sets the interval, and their spread gives the reliability measure. Tests target 90\% coverage.

\q{What are the results?}
Across the 13 real datasets, the spatial version of WBCP reaches a mean coverage of 0.925, essentially tied with BQ-CP (0.928). Its distinguishing output is a reliability measure that varies by location, which BQ-CP cannot give. On the Seattle house-price data, WBCP has the widest intervals (24.9 vs.\ 20.4 for BQ-CP) and its coverage depends on the chosen posterior level (0.873 at $\beta=0.5$, 0.95 at $\beta=0.9$).

\q{What are the impacts and limitations of the contributions?}
WBCP turns a conformal threshold into a distribution under any weighting scheme. The authors state three limitations. Correlated residuals inflate the effective sample size, misspecified weights may void the coverage result, and the theorems are sketches under an assumed model that need a large effective sample size. Yet the method ran with roughly 10--13 effective points on some data.

\q{How were the results evaluated or validated?}
The evaluation uses a synthetic field with known noise, 13 real datasets with random splits, and an ablation. The breadth and the clean ablation are strengths. The weaknesses are no significance tests, a lenient rule for counting a dataset as meeting the target, no check of local coverage, and one figure whose values conflict with the matching table.

\q{What are the ethical implications? How are they addressed?}
There is no ethics statement, the authors discuss neither misuse nor equity, and no code is linked. Because sparsely sampled areas are flagged as less reliable, the reliability map could guide new data collection or be used to justify withholding services from places already poorly measured. The outputs also look like rigorous certificates although they rest on an assumed model and correct weights.

\q{What themes are suggested by this research? (minimum 3)}
\emph{Distribution shift}: calibration scores are reweighted for test inputs that differ from the calibration inputs, though the random splits test no genuine covariate shift. \emph{Uncertainty quantification and calibration}: each threshold reports how well the calibration data support it. \emph{Theory vs.\ practice}: beyond the exchangeable case, WBCP's guarantees rest on an assumed weighted Dirichlet model and a large effective sample size.

\end{qna}
\tiein{WBCP is BCA's calibration procedure, though BCA's code follows the later v3 revision, not the reviewed v2. With BCA's uniform weights it yields one global threshold and is essentially BQ-CP, so the calibrated width inherits an exchangeability assumption that policy-induced shift and correlated transitions can break. Per-sample adaptivity comes from BCA's fitted scale model.}

% ============================ PAPER 13 ============================
\paper{Paper 13: Off-Policy Deep Reinforcement Learning without Exploration (BCQ)}
BCQ shows that standard deep RL algorithms can fail badly on a fixed dataset, and counters this by letting the policy choose only among actions resembling the recorded ones.
\begin{qna}

\q{What is the citation?}
Fujimoto, S., Meger, D., \& Precup, D. (2019). Off-policy deep reinforcement learning without exploration. In \emph{Proceedings of the 36th International Conference on Machine Learning} (Proceedings of Machine Learning Research, Vol.\ 97, pp.\ 2052--2062). PMLR. Version reviewed: arXiv:1812.02900v3. \url{https://doi.org/10.48550/arXiv.1812.02900}

\q{What is the research area?}
The area is offline RL, which the paper calls batch RL (learning from a fixed dataset with no further interaction), for continuous control. BCQ is an early policy-constraint method: it restricts the policy to actions close to those in the data.

\q{What is the research question?}
Why do off-policy deep RL algorithms such as DDPG and DQN fail when trained on a fixed dataset they did not collect themselves? Can restricting the policy to actions similar to the recorded ones make learning from such data work in continuous control?

\q{What is the significance of the question?}
RL needs to learn from a fixed dataset wherever collecting data is costly, risky, or slow. Imitation learning tends to fail on suboptimal demonstrations. Off-policy algorithms seemed to fill that gap, but in practice they train on replay data closely tied to their own current policy, which hides how they behave on data they did not collect.

\q{What are the contributions reported in this paper?}
The paper diagnoses \emph{extrapolation error}: the critic assigns unrealistic values to state-action pairs the dataset lacks, and without new data these are never corrected. Its theory covers finite (tabular), deterministic problems, where a constrained Q-learning (BCQL) converges to the best policy confined to recorded state-action pairs. It presents BCQ as the first deep continuous-control method that learns from arbitrary fixed data, and provides four fixed-data test settings.

\q{What is the research method?}
The paper combines diagnostic experiments, tabular proofs, and benchmarking. BCQ fits a generative model to the recorded actions, samples $n$ candidates, lets a perturbation network shift each by at most $\Phi$, and picks the one the critic scores highest. The tests use three simulated locomotion tasks, with DDPG and behavior cloning (pure imitation) among the baselines.

\q{What are the results?}
DDPG trained on fixed data does markedly worse than the agent that collected it, even when both learn from the same replay buffer. Read off the learning curves, BCQ's return far exceeds the data's average on noisy demonstrations (about 2,300 vs.\ 450 on Hopper, 6,100 vs.\ 2,550 on HalfCheetah). On clean expert data behavior cloning is best on two of three tasks. On two toy tasks with random data, plain DDPG also learns.

\q{What are the impacts and limitations of the contributions?}
The paper names a failure of offline value learning that it calls fundamental, an observation the survey in Literature Review 1 cites. However, the key guarantees cover tabular, deterministic problems, not the deep algorithm, and the authors report instability on one task under slightly different hyperparameters. $\Phi$ is hand-set and, like $n$, the same in every state. The main benchmark data all come from DDPG agents in simulation.

\q{How were the results evaluated or validated?}
The experiments use five seeds and noise-free evaluation, and the learning curves carry half-standard-deviation bands. Value estimates are compared with Monte Carlo returns. There are no result tables or significance tests, and no earlier batch RL method is among the baselines. The appendix shows expert-data returns falling sharply as $\Phi$ grows, but compares $\Phi$ settings by simulator returns, which a purely offline user would lack.

\q{What are the ethical implications? How are they addressed?}
The paper has no ethics or broader-impact discussion. A policy held close to logged behavior may carry over biases and unsafe habits from the log. All evidence is simulated. Public code and the authors' report of instability help, but deployment risks are not discussed.

\q{What themes are suggested by this research? (minimum 3)}
\emph{Distribution shift}: values for actions the data do not contain can be badly wrong and cannot be corrected offline. \emph{The conservatism--improvement trade-off}: $\Phi$ and $n$ slide the method between pure imitation and unconstrained Q-learning. \emph{Theory vs.\ practice}: the key guarantees are proved for tabular, deterministic problems, not for the deep algorithm that is tested.

\end{qna}
\tiein{BCQ is an early form of the ``stay near the data'' idea that two BCA hosts, TD3+BC and ReBRAC, continue as a plain penalty. In BCQ's appendix, critic-ensemble uncertainty did not by itself keep the policy on expert data (one Hopper test), a caution for any method leaning on an uncertainty signal, BCA included.}

% ============================ PAPER 14 ============================
\paper{Paper 14: Stabilizing Off-Policy Q-Learning via Bootstrapping Error Reduction (BEAR)}
BEAR traces much of Q-learning's failure on fixed datasets to value targets at actions the data lack, and keeps the policy approximately within the actions the data support.
\begin{qna}

\q{What is the citation?}
Kumar, A., Fu, J., Soh, M., Tucker, G., \& Levine, S. (2019). Stabilizing off-policy Q-learning via bootstrapping error reduction. \emph{Advances in Neural Information Processing Systems, 32}. Version reviewed: arXiv:1906.00949v2. \url{https://doi.org/10.48550/arXiv.1906.00949}

\q{What is the research area?}
The area is offline RL for continuous control. BEAR belongs to the policy-constraint family with BCQ (Paper 13). Its distinguishing position is support matching (stay within the set of recorded actions) rather than distribution matching (also copy how often each was taken).

\q{What is the research question?}
Why do standard Q-learning and actor-critic methods become unstable on a fixed dataset, even a large or expert one? Can restricting the actions used in value updates to the data's support remove that instability while still letting the policy improve on the data?

\q{What is the significance of the question?}
Logged data are plentiful in areas such as driving and recommendation, where fresh trial and error is costly, yet off-policy methods of the time could fail on a fixed dataset even when an expert collected it. BCQ, an earlier fix, is reported to stay so near the recorded behavior that it gains little over imitation on the mediocre data the authors consider the common practical case.

\q{What are the contributions reported in this paper?}
The paper diagnoses ``bootstrapping error,'' which enters when a value target is computed at an action the data do not cover and accumulates through backups. Its bound for a backup restricted to a chosen set of policies splits the final error into distribution-shift and suboptimality terms. BEAR keeps the actor approximately within the data's support by limiting the maximum mean discrepancy (MMD, a sample-based distance between distributions).

\q{What is the research method?}
The paper combines theory, algorithm design, and benchmarking. Experiments cover five MuJoCo tasks with datasets from a uniformly random policy, a partially trained SAC policy, or a converged one (Humanoid has random and medium data only). Baselines include naive off-policy RL (TD3), BCQ, and behavior cloning.

\q{What are the results?}
On medium-quality data BEAR ends ahead on Walker2d (about 1,800 vs.\ 800 for BCQ), Hopper, and Humanoid, narrowly ahead on HalfCheetah, and just under the dataset's average return on Ant. On random data BEAR beats the dataset average on all five tasks, yet naive RL scores higher on HalfCheetah (about 3,300 vs.\ 2,900) and Hopper. On expert data BEAR, BCQ, and behavior cloning are roughly level.

\q{What are the impacts and limitations of the contributions?}
The paper names the instability and offers a mechanism for it, and CQL (Paper 1) later uses BEAR as a baseline. The authors note that performance can still decay over long runs with no stopping rule, that action constraints can be too conservative, especially on mixed-policy data, and that support matching in high-dimensional action spaces still lacks a theoretically robust method.

\q{How were the results evaluated or validated?}
The theory is proved in an appendix, and the experiments report learning curves averaged over five seeds. There are no tables or significance tests, and ensemble size and other settings are not reported. The paper's claim to match or exceed the state of the art in all cases is not borne out on two random-data tasks. The first term of Theorem 4.1's bound is half what its proof derives.

\q{What are the ethical implications? How are they addressed?}
The paper has no ethics or broader-impact discussion. A policy confined to the data's support can inherit biases of the logged behavior, and the bound concerns value error, not safe behavior. With no offline stopping rule, a policy's quality may be hard to confirm without live trials. All evidence is simulated, and bias, safeguards, and misuse are not addressed.

\q{What themes are suggested by this research? (minimum 3)}
\emph{Distribution shift}: value targets at unseen actions cannot be corrected offline, and backups spread the error. \emph{The conservatism--improvement trade-off}: a wider allowed policy set raises the bound's shift term and lowers its suboptimality term. \emph{Theory vs.\ practice}: the bound assumes policies that truly stay in the support, which a sampled, state-averaged MMD only approximates.

\end{qna}
\tiein{BEAR formalizes, for action constraints, the conservatism--improvement trade-off, which also bears on caution mechanisms such as BCA's. It sets this balance with one threshold, enforced on average over states, whereas BCA varies each host's regularization per sample by a calibrated width. BEAR is not one of BCA's four hosts.}

% ============================ PAPER 15 ============================
\paper{Paper 15: Addressing Extrapolation Error in Deep Offline Reinforcement Learning (BRr)}
BRr learns the value of the policy that recorded the data, so its training targets use only recorded actions, and takes its one improvement step at deployment.
\begin{qna}

\q{What is the citation?}
Gulcehre, C., G\'{o}mez Colmenarejo, S., Wang, Z., Sygnowski, J., Paine, T., \.{Z}o{\l}na, K., Chen, Y., Hoffman, M., Pascanu, R., \& de Freitas, N. (2020). Addressing extrapolation error in deep offline reinforcement learning. In \emph{Offline Reinforcement Learning Workshop at NeurIPS 2020}. Version reviewed: the workshop paper. \url{https://offline-rl-neurips.github.io/pdf/41.pdf} Also available at \url{https://openreview.net/forum?id=OCRKCul3eKN}; a later version appears as \emph{Regularized behavior value estimation} (arXiv:2103.09575).

\q{What is the research area?}
The area is offline RL for problems with a finite set of actions. The paper targets extrapolation error (BCQ's term, Paper 13), and its remedy evaluates only the behavior policy (the policy that recorded the data).

\q{What is the research question?}
Can offline value learning avoid runaway overestimation on poorly covered state-action pairs by estimating the value of the behavior policy and taking a single greedy improvement step only at deployment? With suitable regularization, does this match or beat existing offline methods as data coverage varies?

\q{What is the significance of the question?}
Offline RL is aimed at settings such as healthcare and self-driving cars, where repeated interaction is costly or unsafe. But the max over actions in Q-learning includes actions the data lack, and the critic's guesses for them are sometimes too high. Bootstrapping (building targets from the critic's own predictions) spreads the error, and with no new data nothing corrects it.

\q{What are the contributions reported in this paper?}
The paper contributes three techniques. Behavior value estimation (B) is a SARSA-style update whose target uses the next action actually recorded. The ranking regularizer (R) is a hinge penalty asking the recorded action to outscore every other action. The reparametrized critic (r) is a bounded output times one learned scale. The later arXiv version names the approach R-BVE.

\q{What is the research method?}
The paper combines algorithm design with benchmarking on three discrete-action testbeds. The bsuite tasks have random actions injected at five noise levels. On RL Unplugged Atari, methods are tuned on one set of games and reported on another. DeepMind Lab supplies four levels, plus one level recorded at four noise levels.

\q{What are the results?}
On the held-out Atari games QRr, Q-learning with both regularizers, scores highest (median normalized score 108.2; IQN 104.8). On 10\% of the Atari tuning-game data B scores 58.6 against DDQN's 14.0, though DDQN is slightly ahead on the full data. BRr leads at all five cartpole noise levels. On DeepMind Lab BR's mean return exceeds CQL's on all four levels but falls below behavior cloning (pure imitation) on data recorded with little or no noise.

\q{What are the impacts and limitations of the contributions?}
The paper gives evidence that one greedy step on a behavior-policy critic can compete with dedicated offline methods when coverage is thin. The headline Atari number comes from QRr, which reinstates the max operator; no BRr score is given for the held-out games, and the featured variant changes by domain. BR's return is below each single-level noise dataset's average return, so improvement over the behavior policy is not shown there.

\q{How were the results evaluated or validated?}
The stated protocol sweeps 5--10 hyperparameter settings per method and reports each method's best setting, picked by live evaluation (per level on DeepMind Lab). There are no significance tests, and spreads sometimes overlap (BR 13.4 $\pm$ 11.8 vs.\ CQL 9.3 $\pm$ 2.5). Most Atari baselines are taken from the benchmark paper, and Figure 7 and Table 8 disagree (swapped level labels).

\q{What are the ethical implications? How are they addressed?}
The paper has no ethics or broader-impact section. The deployed policy is greedy over scores that include actions the data never tested, guarded by a heuristic penalty with no stated guarantee. The ranking term favors whatever was done in high-return episodes, so biases in those logs could be reinforced. All evidence is simulated.

\q{What themes are suggested by this research? (minimum 3)}
\emph{Distribution shift}: the behavior-value variants keep it out of the training target by design, because the target bootstraps only from recorded actions. \emph{The conservatism--improvement trade-off}: one step scored higher in the low-data tests but can be too cautious when data are rich. \emph{Theory vs.\ practice}: the motivating domains rule out live trials, yet the best settings were picked by live evaluation.

\end{qna}
\tiein{This workshop paper predates IQL (Paper 3, one of BCA's hosts) and shares its in-sample principle, though the ranking penalty does score unseen actions. Its penalty has one global coefficient, weighted per sample only by trajectory return (data quality, not critic uncertainty), whereas BCA draws its per-sample signal from the size of Bellman residuals.}

% ============================ PAPER 16 ============================
\paper{Paper 16: Uncertainty Weighted Actor-Critic for Offline Reinforcement Learning (UWAC)}
UWAC estimates the critic's uncertainty about each training sample and down-weights uncertain samples in the critic and actor losses, so that value guesses about unseen situations destabilize training less.
\begin{qna}

\q{What is the citation?}
Wu, Y., Zhai, S., Srivastava, N., Susskind, J., Zhang, J., Salakhutdinov, R., \& Goh, H. (2021). Uncertainty weighted actor-critic for offline reinforcement learning. In \emph{Proceedings of the 38th International Conference on Machine Learning} (Proceedings of Machine Learning Research, Vol.\ 139, pp.\ 11319--11328). PMLR. Version reviewed: arXiv:2105.08140v1. \url{https://doi.org/10.48550/arXiv.2105.08140}

\q{What is the research area?}
The area is uncertainty-based offline RL: actor-critic methods that estimate how unreliable the critic is and act on that estimate. UWAC is model-free and builds on BEAR (Paper 14), a policy-constraint method. Its uncertainty tool is Monte Carlo dropout, which leaves dropout's random masking on at prediction time and measures how much repeated predictions disagree.

\q{What is the research question?}
Can a critic's own dropout-based uncertainty identify out-of-distribution (OOD) state-action pairs, meaning ones the dataset does not cover? And does down-weighting those pairs stabilize offline training and raise returns without adding new loss terms or a learned environment model?

\q{What is the significance of the question?}
A value target built from the critic's own prediction at an OOD action or state can be arbitrary. The authors argue that the usual remedy, strong constraints or penalties that hold the policy to the data, avoids unseen actions even when the critic could score them accurately, often leaving policies pessimistic and suboptimal.

\q{What are the contributions reported in this paper?}
The main contribution is a reweighting: each critic update is scaled by the inverse of the estimated variance of its target value, and the actor's objective by the inverse variance of the critic's value for the proposed action. The paper also proposes an explanation of why BEAR's values diverge. The reviewed version links no UWAC code.

\q{What is the research method?}
The method is algorithm design with a short convergence argument, then experiments on D4RL locomotion and Adroit hand-manipulation datasets, with ablations. UWAC adds dropout to BEAR's critics and takes the variance across stochastic forward passes of the target critics as its uncertainty. On Adroit it also adds spectral normalization, a weight constraint the authors use to keep the critic's values bounded.

\q{What are the results?}
Dropout uncertainty separates recorded from random actions on walker2d-expert with an AUC of 0.845 (0.5 is chance). On Adroit, UWAC has the best mean on 8 of 12 datasets and improves on BEAR on human demonstrations (pen-human 65.0 vs.\ $-1.0$), though several human and cloned scores stay near zero. On locomotion it leads on 6 of 15 datasets but trails model-based methods on medium-replay data (hopper 39.4 vs.\ MOReL's 93.6).

\q{What are the impacts and limitations of the contributions?}
Together with spectral normalization, which alone also bounds the critic in the ablation, the uncertainty weighting keeps value targets from diverging on narrow Adroit data. But the dropout variance is never calibrated, only one host algorithm is benchmarked, and as the authors note the convergence argument does not cover the clipped inverse variance actually used. UWAC's pen-human return also decays from about 1,700 to about 700 during training.

\q{How were the results evaluated or validated?}
The results are five-seed D4RL means and standard deviations, with Adroit ablations. Baseline scores appear quoted from published tables, mostly without spreads, and the scored checkpoint is unstated. BEAR is listed at a normalized 98.6 on relocate-expert although Figure 1 shows its average return collapsing to near zero. The text claims state-of-the-art results on every non-human Adroit dataset, yet other methods lead on three cloned ones.

\q{What are the ethical implications? How are they addressed?}
The paper has no ethics or broader-impact statement and no limitations section, and all its evidence is simulated. Down-weighting poorly covered samples lets well-covered behavior shape the policy, so biases in the log could carry over. The uncertainty score could also be read as a safety signal although it is uncalibrated and imperfect by the paper's own test. None of this is discussed.

\q{What themes are suggested by this research? (minimum 3)}
\emph{Distribution shift}: value targets built on unsupported actions are blamed for BEAR's diverging critic on narrow data. \emph{The conservatism--improvement trade-off}: each sample's weight shrinks with its own estimated uncertainty, on top of BEAR's support constraint. \emph{Uncertainty quantification and calibration}: dropout variance is checked as an OOD detector (heat maps, one AUC) but never calibrated or compared with actual value error.

\end{qna}
\tiein{UWAC is an earlier instance of BCA's central move: modulating an existing algorithm's training objective per sample by an uncertainty signal. UWAC's signal is a dropout variance that is never calibrated; BCA instead derives a conformally calibrated width from Bellman residuals. Unlike BCA's current design, UWAC has completed D4RL results.}

% ============================ PAPER 17 ============================
\paper{Paper 17: Uncertainty-Based Offline Reinforcement Learning with Diversified Q-Ensemble (EDAC)}
SAC, an online algorithm, can outscore the tested offline baselines on D4RL locomotion by trusting the lowest of many critics; EDAC's diversity term lets tens replace hundreds on some datasets.
\begin{qna}

\q{What is the citation?}
An, G., Moon, S., Kim, J.-H., \& Song, H. O. (2021). Uncertainty-based offline reinforcement learning with diversified Q-ensemble. \emph{Advances in Neural Information Processing Systems, 34}, 7436--7447. Version reviewed: arXiv:2110.01548v2. \url{https://doi.org/10.48550/arXiv.2110.01548}

\q{What is the research area?}
The area is model-free, uncertainty-based offline RL for continuous control. Caution comes from disagreement among an ensemble of critics instead of tying the policy to the recorded behavior or marking down unseen actions directly. The paper builds on soft actor-critic (SAC) and on clipped Q-learning, the online habit of trusting the lower of two critics.

\q{What is the research question?}
Can disagreement within an ensemble of critics, applied through clipped Q-learning's take-the-minimum rule, supply on its own the pessimism offline RL needs, with no model of the behavior policy and no sampling of unseen actions? And can the number of critics this requires be cut to a practical size?

\q{What is the significance of the question?}
The authors argue that earlier methods typically need an accurate estimate of the behavior policy or samples of out-of-distribution actions. They add that those methods discourage unseen actions wholesale, which can hold the policy too close to mediocre data. A penalty tied to the critics' own confidence could be selective, but the plain ensemble version can demand hundreds of networks.

\q{What are the contributions reported in this paper?}
The paper contributes SAC-$N$, which trains $N$ critics toward a target built from their lowest prediction. It proposes that some datasets need hundreds of critics because, when the critics' slopes with respect to the action line up, nearby unseen actions can escape the penalty. EDAC adds an $\eta$-weighted term that pushes those slopes apart at dataset actions. The reviewed version links no code of its own.

\q{What is the research method?}
The method is algorithm design with a short supporting analysis under simplifying assumptions, then benchmarking. The tests cover 18 D4RL Gym locomotion datasets and eight Adroit robotic-hand datasets, with baselines that include CQL. Diagnostics compare the penalty on dataset vs.\ random actions and find the smallest workable ensemble per dataset.

\q{What are the results?}
On the 18 Gym datasets EDAC averages 85.2 and SAC-$N$ 84.5, against 73.7 for the rerun CQL. On Gym, EDAC's gain over SAC-$N$ is in cost: it uses 50 critics instead of 500 on hopper-medium, where 50 undiversified critics score 7.3. On Adroit, EDAC does well on the pen tasks (pen-cloned: 68.2 vs.\ CQL's published 40.3) but stays near 10 on door and below 1 on hammer and relocate.

\q{What are the impacts and limitations of the contributions?}
The paper shows that on D4RL Gym tasks critic disagreement alone can supply enough pessimism to outscore CQL. But $N$ and $\eta$ are still chosen per environment or dataset, the authors report that a large $\eta$ can hurt on suboptimal data, and the tenfold saving holds on only a few datasets. The spread is not checked against actual value errors, and the analysis gives no performance guarantee.

\q{How were the results evaluated or validated?}
The results are D4RL normalized scores averaged over four seeds. EDAC's Gym settings come from small grids by an unstated criterion, presumably test returns, which a true offline user lacks, and the Adroit stopping point is read off evaluation curves. The rerun CQL is far below its published score on some tasks (hopper-medium: 61.9 vs.\ 86.6).

\q{What are the ethical implications? How are they addressed?}
The reviewed version has no ethics, broader-impact, or limitations section. The method lets the policy leave the data wherever the critics agree, which is arguably a risk. Critics trained on the same data can share an error, so agreement does not show the value is right. All evidence is simulated, and none of these risks is discussed.

\q{What themes are suggested by this research? (minimum 3)}
\emph{The conservatism--improvement trade-off}: the penalty grows with disagreement at each state-action pair instead of treating every unseen action alike. \emph{Uncertainty quantification and calibration}: critic disagreement is the uncertainty signal, and its scale is not checked against real errors. \emph{Theory vs.\ practice}: the reported runs get their pessimism from tens to hundreds of critics, with $N$ and $\eta$ chosen per environment or dataset.

\end{qna}
\tiein{EDAC is not a BCA host but a baseline in the Unifloral codebase bundled with BCA's repository. It represents the ensemble route to per-sample caution through uncalibrated critic spread. BCA seeks comparable selectivity from a conformally calibrated width, with no extra critics, though unlike EDAC its current design has no completed D4RL results yet.}

% ============================ PAPER 18 ============================
\paper{Paper 18: Conformal Prediction as Bayesian Quadrature (BQ-CP)}
BQ-CP rereads conformal prediction as Bayesian numerical integration, so that a calibration set yields a whole distribution bounding a black-box model's expected loss instead of a single point estimate.
\begin{qna}

\q{What is the citation?}
Snell, J. C., \& Griffiths, T. L. (2025). Conformal prediction as Bayesian quadrature. In \emph{Proceedings of the 42nd International Conference on Machine Learning} (Proceedings of Machine Learning Research, Vol.\ 267, pp.\ 56068--56084). PMLR. Version reviewed: arXiv:2502.13228v2. \url{https://doi.org/10.48550/arXiv.2502.13228}

\q{What is the research area?}
The area is statistical machine learning, not RL: distribution-free uncertainty quantification (guarantees on a black-box model's errors that assume no particular data distribution). The paper re-derives split conformal prediction (a threshold set from held-out scores) and conformal risk control (its extension to other bounded, monotone losses) from Bayesian decision theory and Bayesian quadrature.

\q{What is the research question?}
Can split conformal prediction and conformal risk control be recovered as special cases of a Bayesian procedure that needs no hand-picked prior? And can that procedure give guarantees conditioned on the calibration data actually observed, instead of averaged over calibration sets that were never drawn?

\q{What is the significance of the question?}
Black-box predictors are used in high-stakes settings, and a deployment decision rests on the one calibration set in hand. Standard conformal guarantees are marginal: they hold on average over hypothetical calibration sets, so a particular set can still produce a threshold whose real risk is too high.

\q{What are the contributions reported in this paper?}
The paper frames both conformal methods decision-theoretically, as keeping worst-case risk below a target $\alpha$. Using a classical result that gaps between the quantile levels of i.i.d.\ samples follow a flat Dirichlet distribution, it defines a random variable $L^+$ whose distribution upper-bounds the posterior of the expected loss. Both conformal methods are recovered by using only its mean.

\q{What is the research method?}
The method is theory followed by simulation. The proofs yield a decision rule: choose the smallest threshold $\lambda$ for which $L^+$ stays below the target risk $\alpha$ with probability $\beta = 0.95$. The rule is tested over 10,000 random splits in each of three settings: a synthetic binomial loss, a synthetic regression, and false-negative-rate control on MS-COCO. The baselines are conformal risk control (CRC) and risk-controlling prediction sets (RCPS).

\q{What are the results?}
In theory, given the calibration losses, the expected loss is at most $b$ with posterior probability at least $\Pr(L^+ \le b)$ under any prior compatible with them. Empirically, CRC's threshold exceeds the target risk in 21.20\%, 46.19\%, and 45.05\% of trials across the three settings; BQ-CP cuts this to 0.03\%, 3.42\%, and 5.43\%, the last just above its 5\% allowance. Its intervals are wider than CRC's (mean length 9.50 vs.\ 7.99; RCPS: 14.29).

\q{What are the impacts and limitations of the contributions?}
BQ-CP received an outstanding paper award at ICML 2025 and is the base WBCP (Paper 12) extends. The authors state that deployment data must be i.i.d.\ with the calibration data and that bounds valid for any prior can be loose. Beyond these, the rule returns one global threshold with no per-input guarantee, and no theorem bounds how often it fails over repeated calibration sets.

\q{How were the results evaluated or validated?}
The results are validated by proofs and repeated sampling. The proofs assume continuous losses; discrete losses, such as those of the two synthetic experiments, get only an unproved footnote. One step appears to take the Dirichlet law of the gaps, a sampling result, as their posterior once the losses are seen, without arguing this for every prior. The experiments use one real dataset and a single $\beta$.

\q{What are the ethical implications? How are they addressed?}
A short impact statement says that the method's guarantees may fail if its assumptions are violated and that deployments should be monitored for distribution shift. Beyond the statement, a posterior bound can be mistaken for a safety certificate when the i.i.d.\ assumption goes unchecked, and a bound on average loss can hide high loss for a subgroup. The statement offers no monitoring method or subgroup analysis.

\q{What themes are suggested by this research? (minimum 3)}
\emph{The conservatism--improvement trade-off}: RCPS records no failures but is widest, CRC is tightest but often exceeds the target, and BQ-CP sits between. \emph{Uncertainty quantification and calibration}: a calibration set yields a distribution that upper-bounds the posterior of the expected loss. \emph{Theory vs.\ practice}: CRC's marginal guarantee can hold while, in the paper's three experiments, CRC's threshold exceeds the target risk for 21--46\% of calibration sets.

\end{qna}
\tiein{BQ-CP is in effect the threshold rule BCA runs, at this paper's $\beta = 0.95$. BCA inherits the guarantee only where the i.i.d.\ assumption holds, and it is neither a per-sample nor, as BCA's own documentation says, a policy-performance guarantee. Same-episode, policy-dependent Bellman residuals strain that assumption even after BCA thins its calibration set.}

\section*{Part A: Literature Reviews}

The three literature reviews are new in Version 2. Each is reviewed against the eleven literature-review questions, and their themes and meta-themes are drawn from the same four themes.

% ====================== LITERATURE REVIEW 1 ======================
\paper{Literature Review 1: Offline Reinforcement Learning: Tutorial, Review, and Perspectives on Open Problems (Levine et al., 2020)}
This early tutorial explains why learning from a fixed dataset is hard, groups early methods by how they cope, and lists open problems.
\begin{qna}

\q{What is the citation?}
Levine, S., Kumar, A., Tucker, G., \& Fu, J. (2020). \emph{Offline reinforcement learning: Tutorial, review, and perspectives on open problems} (arXiv:2005.01643v3) [Preprint]. arXiv. \url{https://doi.org/10.48550/arXiv.2005.01643}

\q{What is the survey scope?}
The survey covers offline RL as a whole: learning a policy from a fixed dataset with no interaction during training. It reviews three method families: importance-sampling estimators; approximate dynamic programming (critic-based value learning) with policy constraints, uncertainty penalties, or conservative critics; and model-based methods. It also covers evaluation practice and applications, cites roughly 190 references, and sets aside exploration and a complete account of model-based RL.

\q{What is the survey question?}
What prevents current RL algorithms from turning a fixed, previously collected dataset into the best policy the data can support, and which ideas have been proposed to overcome those obstacles?

\q{What is the significance of the survey?}
It offered newcomers one entry point, placing older statistical work on importance sampling beside recent deep methods under one notation. It appeared in May 2020, when offline RL had just become an active topic (BCQ and BEAR, Papers 13 and 14). Figueiredo Prudencio et al.\ later treat it as their main predecessor.

\q{What are the themes?}
\emph{Distribution shift}: this is the core obstacle, since learning offline means answering what-if questions about actions the data never contained. \emph{The conservatism--improvement trade-off}: policy constraints keep the learned policy near the data, and conservative critics make the critic itself cautious. \emph{Uncertainty quantification and calibration}: uncertainty-based methods make the critic cautious through an explicit uncertainty model.

\q{What is the survey method?}
It is an expert narrative written as a tutorial, not a systematic review: it reports no search databases, keywords, or inclusion criteria. The authors restate and explain results from the literature but prove no new theorems and run no experiments. Their own critical view appears mainly in the challenges subsections.

\q{What are the survey findings?}
Importance-sampling variance grows exponentially with the horizon, so those tools are most suitable when the learned policy stays close to the data. Naive dynamic programming can perform very poorly because its targets query unseen actions. Calibrated uncertainty for deep value functions is hard to obtain, so constraint and conservative methods seemed to outperform uncertainty-based ones at the time. Conservative critics can be too pessimistic. The survey judges benchmarking immature.

\q{What are the limitations of the survey?}
It stops at mid-2020, before TD3+BC, IQL, ReBRAC, and the sequence-model methods. Coverage is uneven, with uncertainty-based methods given about half a page. More than 30 of the roughly 190 references include one of the four authors, whose own methods (BEAR, BRAC, CQL) are among the families the survey reports as performing best. It offers no quantitative comparison and never treats offline hyperparameter selection directly.

\q{How was the survey assessed relative to other surveys?}
It has no related-surveys section. Of the three surveys reviewed here, it gives the most formal and conceptual account of \emph{why} offline RL is hard. The survey by Figueiredo Prudencio et al.\ (2024) is broader and more systematic, with a taxonomy and a benchmark review. He's survey (2023) is a short list of model-based methods with little analysis.

\q{What are the ethical implications? How are they addressed?}
The survey has no ethics section. Safety appears only in scattered technical remarks, as a motivation rather than an ethical question. The survey does not address privacy and consent for clinical data, unequal treatment of groups under-represented in logs, or the danger of treating off-policy estimates as the final check before deployment.

\q{What meta-themes are suggested by this survey? (minimum 3)}
\emph{Distribution shift}: the survey traces each family's main difficulty to this mismatch, so later methods can be compared by where they intercept it. \emph{The conservatism--improvement trade-off}: tight constraints are needed to keep errors small, yet they tie the policy to possibly poor behavior. \emph{Uncertainty quantification and calibration}: this is a principled remedy that is hard to deliver, and dynamic modulation of conservative critics is left open.

\end{qna}
\tiein{The survey lays out the method families of BCA's four host algorithms. BCA takes up the survey's open problem of how to dynamically modulate conservatism. Where the survey judged uncertainty estimates hard to calibrate, BCA uses residual widths calibrated on held-out transitions, though that calibration covers recorded transitions, not unseen actions.}

% ====================== LITERATURE REVIEW 2 ======================
\paper{Literature Review 2: A Survey on Offline Model-Based Reinforcement Learning (He, 2023)}
This short primer on offline model-based RL profiles seven methods by how each builds caution into learning with a model of the environment.
\begin{qna}

\q{What is the citation?}
He, H. (2023). \emph{A survey on offline model-based reinforcement learning} (arXiv:2305.03360v1) [Preprint]. arXiv. \url{https://doi.org/10.48550/arXiv.2305.03360}

\q{What is the survey scope?}
The survey covers offline model-based RL: fitting a dynamics and reward model to a fixed log of experience, then planning or optimizing a policy inside that model. It profiles seven methods from 2020 to early 2023 (MOReL, MOPO, BREMEN, VI-LCB, COMBO, RAMBO, ARMOR) in about five pages, with nineteen references. Benchmarks, off-policy evaluation, and application domains are absent.

\q{What is the survey question?}
How have recent offline model-based methods coped with the distribution shift that arises when a model learned from logged data is used to optimize a new policy, and what problems remain unsolved?

\q{What is the significance of the survey?}
He's paper is a dedicated, recent, and very short treatment of offline model-based RL, which the two general surveys reviewed here gave only a few pages. It is a single-author preprint with no peer-reviewed version listed on arXiv, so it serves as a quick map rather than an authoritative reference.

\q{What are the themes?}
\emph{Distribution shift}: a policy optimized inside a model learned from logged data may be out of distribution at deployment and do worse. \emph{The conservatism--improvement trade-off}: COMBO's conservative critic and the adversarially trained models of RAMBO and ARMOR supply caution by means other than uncertainty penalties. \emph{Uncertainty quantification and calibration}: MOReL and MOPO turn ensemble uncertainty into a reward penalty.

\q{What is the survey method?}
It is a personal, narrative selection: no databases, search terms, or inclusion criteria are named, and the choice of these seven methods is not explained. Each profile gives the core idea, the key equations in the notation of whichever source it follows, and the original authors' claims. There are no experiments or new results. Some passages follow Figueiredo Prudencio et al.\ closely.

\q{What are the survey findings?}
By the survey's account, distribution shift, made worse by model exploitation, remains the defining open problem, and none of the reviewed methods provides a performance guarantee that holds regardless of the data. Reward-penalty methods depend on uncertainty estimates the survey regards as error-prone. Citing the ARMOR paper, it adds that many theoretically grounded methods go unused because they fail to beat the policy already in service.

\q{What are the limitations of the survey?}
Coverage is thin, with no benchmark results or comparison table. Most profiles offer little independent critique. Some descriptions are loose: COMBO is said to penalize rewards, while the objective the survey prints acts on Q-values. It has not been revised since its first version.

\q{How was the survey assessed relative to other surveys?}
It cites the other two surveys only as background authorities and never says what it adds to them. Levine et al.\ give model-based methods a more analytical section, and Figueiredo Prudencio et al.\ place them in a full taxonomy. He's distinct additions are RAMBO, ARMOR, VI-LCB, and the tabular sample-complexity result, none of which appears in either broader survey.

\q{What are the ethical implications? How are they addressed?}
The survey discusses none. It frames offline learning as improving on a deployed system. There, model exploitation means a policy can look excellent inside the learned simulator and fail once it replaces that system. Biases in past decisions can pass into both the model and the new policy. The survey never explains how such policies could be checked before release; only the safe-improvement guarantees of COMBO and ARMOR come near the topic.

\q{What meta-themes are suggested by this survey? (minimum 3)}
\emph{The conservatism--improvement trade-off}: several methods vary caution across state-action pairs, but where the survey says how its strength is set, it is hand-chosen. \emph{Uncertainty quantification and calibration}: MOReL's and MOPO's penalties rest on ensemble heuristics whose link to actual model error the survey does not examine. \emph{Theory vs.\ practice}: the only optimality result the survey reports is tabular.

\end{qna}
\tiein{He's survey suggests the model-based branch, like BCA's model-free hosts, sets its caution by hand-chosen hyperparameters and sizes some penalties by uncertainty estimates it calls possibly inaccurate. Because a method like COMBO mirrors CQL's penalty, BCA's per-sample widths might carry over to this branch, which is untested.}

% ====================== LITERATURE REVIEW 3 ======================
\paper{Literature Review 3: A Survey on Offline Reinforcement Learning: Taxonomy, Review, and Open Problems (Figueiredo Prudencio et al., 2024)}
This broad, structured update to Levine et al.\ offers a taxonomy of offline algorithms, a review of benchmarks, and a refreshed list of open problems.
\begin{qna}

\q{What is the citation?}
Figueiredo Prudencio, R., Maximo, M. R. O. A., \& Colombini, E. L. (2024). A survey on offline reinforcement learning: Taxonomy, review, and open problems. \emph{IEEE Transactions on Neural Networks and Learning Systems, 35}(8), 10237--10257. \url{https://doi.org/10.1109/TNNLS.2023.3250269}

\q{What is the survey scope?}
The survey covers deep offline RL in the single-agent setting: policy constraints, importance sampling, regularization, uncertainty estimation, model-based methods, one-step methods, imitation learning, and trajectory optimization. Its timeline runs from late 2018 to early 2022, its method table classifies 26 entries, and it cites 91 references. It also reviews off-policy evaluation and the D4RL and RL Unplugged benchmarks. It leaves out theoretical analyses.

\q{What is the survey question?}
How can the fast-growing body of offline RL algorithms be placed in one framework that shows each method's design choices, how those choices fare on different kinds of data, and where research effort should go next?

\q{What is the significance of the survey?}
The survey supplies a taxonomy that lets methods be read as combinations of reusable parts. It also maps dataset properties to the algorithm classes that handle them, which is meant to help practitioners match methods to their data.

\q{What are the themes?}
\emph{Distribution shift}: most classes are introduced as ways to prevent, penalize, or sidestep this central difficulty. \emph{The conservatism--improvement trade-off}: this appears as policy constraints and behavior regularization, with direct constraints set against implicit ones that use only dataset samples. \emph{Theory vs.\ practice}: this concerns evaluating policies without deployment, with off-policy evaluation tied to hyperparameter selection and early stopping.

\q{What is the survey method?}
It is an expert survey that names no databases, search strings, or inclusion criteria, and the authors note that their method table is far from exhaustive. Methods are restated in a shared notation. The only quantitative analysis is secondary: a heatmap built from D4RL scores collected from the literature. No new experiments are run.

\q{What are the survey findings?}
Most new algorithms are different responses to distribution shift, each with a trade-off. Policy constraints can be overly pessimistic, and direct ones also need a good estimate of the behavior policy. The survey relays two lessons: hyperparameter search mattered more than the choice of divergence (from BRAC), and performance is often not tied to algorithmic complexity (from TD3+BC).

\q{What are the limitations of the survey?}
The uncertainty section reviews only one line of ensemble work. The heatmap pools scores from many papers, so it likely mixes tuning budgets and codebases; it also credits each class with its best member. Methods after early 2022, such as ReBRAC, are absent. At least 21 of the 91 references share a single coauthor from another research group, which may tilt the review toward that group.

\q{How was the survey assessed relative to other surveys?}
It positions itself against one predecessor, Levine et al.\ (2020), whose tutorial remains longer and more conceptual. Of the three surveys reviewed here it gives the broadest map of model-free methods and is the only one to systematically map benchmark environments to dataset properties. He's preprint is far narrower but reaches theory-driven and adversarial model-based work that this survey omits.

\q{What are the ethical implications? How are they addressed?}
The survey has no ethics section. Safety appears as a motivation and as an open problem (safety-critical RL). The survey asks for datasets with risky gaps to test whether methods learn to avoid hazards missing from the data. It does not discuss prejudice carried in logged human decisions, privacy, consent, or accountability. As scholarship it discloses funding and avoids self-citation.

\q{What meta-themes are suggested by this survey? (minimum 3)}
\emph{The conservatism--improvement trade-off}: in its equations most mechanisms carry one scalar weight or threshold, and dynamic modulation of conservatism is named but not developed. \emph{Uncertainty quantification and calibration}: the survey calls model uncertainty estimates often unreliable but does not discuss calibration (whether an estimate matches observed error). \emph{Theory vs.\ practice}: off-policy evaluation and hyperparameter selection lead the list of open problems.

\end{qna}
\tiein{This survey explains three of BCA's four host algorithms side by side (CQL, TD3+BC, IQL). It names, without developing it, the open problem BCA targets: dynamically modulating the degree of conservatism. Its reliance on D4RL is also a caveat for BCA, since any results measured on that benchmark alone carry its limits.}

\section*{Part B: Themes and Observations}

All twenty-one works (eighteen research papers and three literature reviews) are described with the same four themes. Each paper names three, and for a review I count its themes and meta-themes together.

\emph{Theory vs.\ practice} is the most common (19 works). What is proved or reported often differs from what an offline user gets: in several papers the guarantees assume conditions the tested algorithm does not meet, and scores depend on implementation details and live tuning.

\emph{The conservatism--improvement trade-off} comes next (18 works). An offline method has to decide how far to distrust values its data cannot verify, and many methods set that caution with one hand-chosen number.

\emph{Distribution shift} (16 works) is the problem the other themes respond to: values for actions the data do not contain cannot be corrected offline.

\emph{Uncertainty quantification and calibration} is the least common theme (12 works) but the most relevant to this project, because it offers a way to make caution vary per sample. The three surveys describe uncertainty estimates as hard to obtain or often unreliable. Five works add a calibration step (Papers 8, 10, 11, 12, and 18), and this review notes limits on each guarantee.

These themes motivate BCA and also bound it. BCA is designed to set caution per sample using calibrated uncertainty, but a calibration guarantee is not a policy-performance guarantee, and BCA's benefit has not yet been shown on benchmarks.

\end{document}
